\section{Related Work}
\label{sec:related}

\paragraph{Orthogonalized updates.}
Muon replaces the momentum matrix by its polar factor before the weight
update~\cite{liu2025muon}; that factor is computed by a fixed Newton--Schulz-type
iteration whose convergence and error analysis are
classical~\cite{higham1986polar,higham2008functions,golub2013matrix}, and
Shampoo~\cite{gupta2018shampoo} is the other line of preconditioned tensor optimization.
Amsel et al.~\cite{amsel2025polar} derive minimax quintic coefficients for exact
arithmetic and add the prescriptions this paper takes as its object; the empty design
axis (Appendix~\ref{sec:design}) is consistent with that optimality holding only in
exact arithmetic.

\paragraph{Circuits against the chain, and two different certificates.}
Du and Su~\cite{du2026polarswitch} replace the quintic chain by polynomial circuits of
degree up to $65$, choosing the products and coefficients to minimize a worst-case error
under a FLOP budget and a half-precision rounding constraint; they report $15.5$--$18.7\%$
less runtime at comparable accuracy, with the output's spectral norm inside its designed
bound on $14$ million matrices.  The certificates differ: a circuit buys FLOPs at a fixed
singular-value error, whereas the subspace error is no functional of the singular values
(Proposition~\ref{prop:nospectral}) and no polynomial removes it at fixed bits
(Corollary~\ref{cor:flat}).  Our numbers say where such a design stops paying: the shipped
point is within $1.11\times$ of the best of 35 designs at bf16, a step past $K=5$ costs
eight bits for under $1.5\%$, and the four-bit failure is an escape from the window rather
than an accuracy loss (\S\ref{sec:escape}).
Huang~\cite{huang2026orthogonalization} reports a non-monotone relation between
orthogonalization accuracy and training quality, which the two-term law explains in shape
if not in constant.

\paragraph{Low precision and bit allocation.}
The formats are standardized~\cite{micikevicius2022fp8} and mixed-precision training is
standard~\cite{micikevicius2018mixed}; optimizer states have been quantized
blockwise~\cite{dettmers2022optimizers}, and limited precision with stochastic rounding
has been studied for training~\cite{gupta2015limited}.  That line asks what precision the \emph{parameters and states} need; we ask what the
\emph{working matrices of one iteration} need, and the answer is a floor proportional to $u$
with a computable constant and a separate exponent-range test.  Quantization work
allocates bits across layers, channels or tensors~\cite{hariri2025quantize,dettmers2022optimizers};
here allocation is inside one step, across $K+1$ slots, and it is a maximum rather than a
sum, which makes a closed form possible.

\paragraph{Perturbation of the polar factor, and stopping.}
First-order perturbation of the SVD and of orthonormal bases is
classical~\cite{stewart1990perturbation,golub2013matrix}; the two-term form
\eqref{eq:firstorder}, with the projected term written separately because it survives for
rectangular matrices, is what we could not find stated as a working estimate for a
rounded momentum matrix.  Early-exit rules in iterative solvers test a residual the algorithm can see; here the law
runs before the iteration, against an error it cannot see.
